% part: intuitionistic-logic
% chapter: introduction
% section: constructive-reasoning

\documentclass[../../../include/open-logic-chapter]{subfiles}

\begin{document}

\olfileid[hi]{int}{int}{cr}

\section{रचनात्मक तर्क-विचार}

मोडल संकारकों या द्वितीय-कोटि परिमाणकों द्वारा शास्त्रीय तर्कशास्त्र का
विस्तार करने के विपरीत, अंतःप्रज्ञावादी तर्कशास्त्र इस अर्थ में
``गैर-शास्त्रीय'' है कि वह शास्त्रीय तर्कशास्त्र को सीमित करता है। शास्त्रीय
तर्कशास्त्र अनेक प्रकार से \emph{अरचनात्मक} है। अंतःप्रज्ञावादी तर्कशास्त्र
का उद्देश्य उस अधिक ``रचनात्मक'' तर्क-विचार को ग्रहण करना है जो एक प्रकार
के रचनात्मक गणित की विशेषता है। निम्न उदाहरण इसके कुछ आधारभूत प्रेरकों को
स्पष्ट कर सकते हैं।

मान लें कि कोई व्यक्ति दावा करता है कि उसने ऐसा प्राकृतिक संख्या~$n$
निर्धारित कर लिया है कि यदि $n$ सम है तो रीमान परिकल्पना सत्य है, और यदि
$n$ विषम है तो रीमान परिकल्पना असत्य है। यह तो बड़ी अच्छी खबर है!{} रीमान
परिकल्पना सत्य है या नहीं, गणित के बड़े खुले प्रश्नों में से एक है; और लगता
है कि उसने समस्या को मात्र एक गणना में बदल दिया है---अर्थात यह निर्धारित
करने में कि कोई विशेष संख्या सम है या नहीं।

तो~$n$ का वह जादुई मान क्या है? वह इसे इस प्रकार बताता है: $n$ वह प्राकृतिक
संख्या है जो रीमान परिकल्पना सत्य होने पर $2$, और अन्यथा $3$ के बराबर है।

क्रोधित होकर आप अपने पैसे वापस माँगते हैं। शास्त्रीय दृष्टिकोण से ऊपर का
वर्णन वास्तव में $n$ का एक अद्वितीय मान निर्धारित करता है; किंतु आप सचमुच
$n$ का ऐसा मान चाहते हैं जो \emph{स्पष्टतः} दिया गया हो।

एक और, शायद कम कृत्रिम, उदाहरण के लिए निम्न प्रश्न पर विचार करें। हम जानते
हैं कि किसी अपरिमेय संख्या को परिमेय घात तक उठाकर परिमेय परिणाम पाना संभव
है; उदाहरणार्थ $\sqrt{2}^2 = 2$। यह कम स्पष्ट है कि क्या किसी अपरिमेय
संख्या को \emph{अपरिमेय} घात तक उठाकर परिमेय परिणाम पाना संभव है। निम्न
प्रमेय इसका उत्तर हाँ में देता है:

\begin{thm}
ऐसी अपरिमेय संख्याएँ $a$ और $b$ हैं कि $a^b$ परिमेय है।
\end{thm}

\begin{proof}
$\sqrt{2}^{\sqrt{2}}$ पर विचार करें। यदि यह परिमेय है, तो काम पूरा हुआ:
हम $a = b = \sqrt{2}$ ले सकते हैं। अन्यथा यह अपरिमेय है। तब हमारे पास
\[
(\sqrt{2}^{\sqrt{2}})^{\sqrt{2}} = \sqrt{2}^{\sqrt{2} \cdot
  \sqrt{2}} = \sqrt{2}^2 = 2,
\]
है, जो परिमेय है। अतः इस स्थिति में $a$ को $\sqrt{2}^{\sqrt{2}}$ और $b$
को~$\sqrt 2$ लें।
\end{proof}

क्या यह एक वैध प्रमाण है? अधिकांश गणितज्ञ मानते हैं कि हाँ। किंतु यहाँ फिर
कुछ असंतोषजनक है: हमने किसी निश्चित गुण वाले वास्तविक संख्याओं के युग्म का
अस्तित्व सिद्ध किया, पर यह नहीं बता सके कि वह \emph{कौन-सा} युग्म है। उसी
परिणाम को ऐसे भी सिद्ध किया जा सकता है कि प्रमाण में युग्म $a$, $b$
\emph{दिया} जाए: $a = \sqrt{3}$ और $b = \log_3 4$ लें। तब
\[
a^b = \sqrt{3}^{\log_3 4} = 3^{1/2 \cdot \log_3 4} = (3^{\log_3
  4})^{1/2} = 4^{1/2}= 2,
\]
क्योंकि $3^{\log_3 x} = x$।

अंतःप्रज्ञावादी तर्कशास्त्र उस प्रकार के तर्क-विचार को ग्रहण करने के लिए
बनाया गया है जिसमें पहले प्रमाण जैसी चालों की अनुमति नहीं होती। $!A(x)$ को
संतुष्ट करने वाले किसी~$x$ का अस्तित्व सिद्ध करने का अर्थ है कि आपको कोई
विशिष्ट~$x$ और उसके द्वारा $!A$ को संतुष्ट करने का प्रमाण देना होगा, जैसा
दूसरे प्रमाण में हुआ। $!A$ या $!B$ सत्य है, यह सिद्ध करने के लिए आपको उन
दोनों में से किसी एक को सिद्ध कर सकना होगा।

औपचारिक रूप से कहें तो अंतःप्रज्ञावादी तर्कशास्त्र, शास्त्रीय तर्कशास्त्र की
किसी !!{derivation}-प्रणाली को एक निश्चित प्रकार से सीमित करने पर मिलता
है। गणितीय दृष्टि से ये केवल औपचारिक निगमनात्मक प्रणालियाँ हैं, किंतु जैसा
पहले कहा गया, उनका उद्देश्य एक प्रकार के गणितीय तर्क-विचार को ग्रहण करना
है। इसे हम गणित के किसी दार्शनिक दृष्टिकोण (जैसे ब्राउवर के अंतःप्रज्ञावाद)
पर उचित ठहरने वाला तर्क-विचार मान सकते हैं; इसे अधिक ``मूर्त'' और संतोषजनक
गणितीय तर्क-विचार (बिशप के रचनावाद की दिशा में) मान सकते हैं; और इस पर
बहस कर सकते हैं कि औपचारिक वर्णन अनौपचारिक प्रेरणा को ग्रहण करता है या नहीं।
किंतु हमारे दार्शनिक मत जो भी हों, हम अंतःप्रज्ञावादी तर्कशास्त्र का अध्ययन
औपचारिक रूप से प्रस्तुत तर्कशास्त्र के रूप में कर सकते हैं; और कारण जो भी
हों, अनेक गणितीय तर्कशास्त्री ऐसा करना रोचक मानते हैं।

\end{document}
